% !PRINTED 104
\weberpage{104}
%<<CONT
are indicated, and whose number one finds, just as above, to be
\[
= \frac{3^{\mu - 2} - 1}{2}
\]
If, then, $m$ is divisible by 8, the total number of the groups $\mathfrak{A}$ is
\weq{5}{\frac{3^{\mu - 1} - 1}{2} + \frac{3^{\mu - 2} - 1}{2} = 2 . 3^{\mu - 2} - 1.}
The smallest values that $m$ can assume in this case are
\[
m = 8, 12, 15, 20, 21, 24, 28, 33, 35, 39, 40, 44 \ldots
\]
The smallest number that gives rise to more than one such group is $m = 24$, for which one has $\mu = 3$, and which therefore, by (5), gives five different groups. One obtains, namely, the following admissible and mutually distinct combinations of the $e_{h,k}$: % [sic: Weber again prints e_{h,k}; the displayed combinations and defining congruences use e_{k,h}.]
\[
\begin{array}{ccccc}
1\ 0\ 1 & 1\ 1\ 0 & 1\ 1\ 0 & 1\ 0\ 0 & 1\ 0\ 0\\
0\ 1\ 0, & 0\ 0\ 1, & 0\ 1\ 1, & 0\ 1\ 1, & 0\ 0\ 1,
\end{array}
\]
which leads to the following determinations concerning the indices:
\[
\begin{aligned}
& \alpha_1 \equiv \alpha_3,\quad \alpha_2 \equiv 0,\\
& \alpha_1 \equiv \alpha_2,\quad \alpha_3 \equiv 0,\\
& \alpha_1 \equiv \alpha_2 \equiv \alpha_3,\qquad (\mathrm{mod}\ 2)\\
& \alpha_1 \equiv 0,\quad \alpha_2 \equiv \alpha_3\\
& \alpha_1 \equiv 0,\quad \alpha_3 \equiv 0,
\end{aligned}
\]
and if one therefore determines the numbers $a$ by the congruences
\[
a \equiv (-1)^{\alpha_2} 5^{\alpha_1} (\mathrm{mod}\ 8), \quad a \equiv 2^{\alpha_3} (\mathrm{mod}\ 3)
\]
the following five two-term groups result:
\[
\begin{aligned}
\mathfrak{A}_1 &= 1,\ 5\\
\mathfrak{A}_2 &= 1,\ 19\\
\mathfrak{A}_3 &= 1,\ 11\\
\mathfrak{A}_4 &= 1,\ 23\\
\mathfrak{A}_5 &= 1,\ 7.
\end{aligned}
\]
To compute the periods $\eta$, we can proceed similarly as in the preceding section. We denote by $\nu_1, \nu_2, \ldots \nu_\mu$ the indices of a number $n$ and set
\[
\begin{aligned}
e_{1,1} \nu_1 + e_{1,2} \nu_2 + \cdots + e_{1,\mu} \nu_\mu &= \sigma_1,\\
e_{2,1} \nu_1 + e_{2,2} \nu_2 + \cdots + e_{2,\mu} \nu_\mu &= \sigma_2,
\end{aligned}
\]
%CONT>>
